% GENERATED FILE. Edit canonical lessons, then run npm run generate:books.
% canonical-source-hash: fnv1a64:4e369850319794fe
% canonical-lessons: KA-C182-jipuna, KA-C182-hathamari, KA-C182-koleta, KA-C182-magida, KA-C182-svanta

\chapter{Stingy, Stubborn, Rotten, Ripe, One's own}
\label{ch:ka-stingy-stubborn-rotten-ripe-one-s-own}
\hlchaptermodality{182}

\begin{chapteropening}
\textbf{By the end of this chapter:} \emph{I can say stingy, stubborn, rotten, ripe and one's own.}

Complete the last lesson of chapter 182: I can say stingy, stubborn, rotten, ripe and one's own.
\end{chapteropening}

\section[\texorpdfstring{\kn{ಜಿಪುಣ} (jipuṇa)}{ಜಿಪುಣ (jipuṇa)}]{\kn{ಜಿಪುಣ} (jipuṇa) — stingy, miserly}

\label{lesson:KA-C182-jipuna}



\begin{quote}
Before the new one: say the Kannada for modern, then the Kannada for traditional.
\end{quote}

\subsection*{You'll want to know: \kn{ಜಿಪುಣ}}
\textbf{\kn{ಜಿಪುಣ}} — \emph{jipuṇa} — \textquotedblleft{}stingy, miserly\textquotedblright{}.

\begin{grammarlens}[title={What you've built}]
One more word for this chapter.
\end{grammarlens}

\subsection*{Your turn}
\begin{itemize}
\raggedright
  \item \emph{Say it:} \emph{jipuṇa}
  \item \emph{Say it:} \emph{jipuṇa}, once more
  \item \emph{Say it:} say \emph{jipuṇa}
  \item \emph{Recall:} say the Kannada for modern, then the Kannada for traditional, then say \emph{jipuṇa} again
\end{itemize}

\begin{tcolorbox}[breakable,colback=teal!4,colframe=teal!35!black,title={Before you move on}]
What does \kn{ಜಿಪುಣ} mean? (\textquotedblleft{}Stingy, miserly\textquotedblright{}.) Say it once more.
\end{tcolorbox}
\section[\texorpdfstring{\kn{ಹಠಮಾರಿ} (haṭhamāri)}{ಹಠಮಾರಿ (haṭhamāri)}]{\kn{ಹಠಮಾರಿ} (haṭhamāri) — stubborn}

\label{lesson:KA-C182-hathamari}



\begin{quote}
Before the new one: say the Kannada for traditional, then the Kannada for stingy.
\end{quote}

\subsection*{You'll want to know: \kn{ಹಠಮಾರಿ}}
\textbf{\kn{ಹಠಮಾರಿ}} — \emph{haṭhamāri} — \textquotedblleft{}stubborn\textquotedblright{}.

\begin{grammarlens}[title={What you've built}]
One more word for this chapter.
\end{grammarlens}

\subsection*{Your turn}
\begin{itemize}
\raggedright
  \item \emph{Say it:} \emph{haṭhamāri}
  \item \emph{Say it:} \emph{haṭhamāri}, once more
  \item \emph{Say it:} say \emph{haṭhamāri}
  \item \emph{Recall:} say the Kannada for traditional, then the Kannada for stingy, then say \emph{haṭhamāri} again
\end{itemize}

\begin{tcolorbox}[breakable,colback=teal!4,colframe=teal!35!black,title={Before you move on}]
What does \kn{ಹಠಮಾರಿ} mean? (\textquotedblleft{}Stubborn\textquotedblright{}.) Say it once more.
\end{tcolorbox}
\section[\texorpdfstring{\kn{ಕೊಳೆತ} (koḷeta)}{ಕೊಳೆತ (koḷeta)}]{\kn{ಕೊಳೆತ} (koḷeta) — rotten}

\label{lesson:KA-C182-koleta}



\begin{quote}
Before the new one: say the Kannada for stingy, then the Kannada for stubborn.
\end{quote}

\subsection*{You'll want to know: \kn{ಕೊಳೆತ}}
\textbf{\kn{ಕೊಳೆತ}} — \emph{koḷeta} — \textquotedblleft{}rotten\textquotedblright{}.

\begin{grammarlens}[title={What you've built}]
One more word for this chapter.
\end{grammarlens}

\subsection*{Your turn}
\begin{itemize}
\raggedright
  \item \emph{Say it:} \emph{koḷeta}
  \item \emph{Say it:} \emph{koḷeta}, once more
  \item \emph{Say it:} say \emph{koḷeta}
  \item \emph{Recall:} say the Kannada for stingy, then the Kannada for stubborn, then say \emph{koḷeta} again
\end{itemize}

\begin{tcolorbox}[breakable,colback=teal!4,colframe=teal!35!black,title={Before you move on}]
What does \kn{ಕೊಳೆತ} mean? (\textquotedblleft{}Rotten\textquotedblright{}.) Say it once more.
\end{tcolorbox}
\section[\texorpdfstring{\kn{ಮಾಗಿದ} (māgida)}{ಮಾಗಿದ (māgida)}]{\kn{ಮಾಗಿದ} (māgida) — ripe}

\label{lesson:KA-C182-magida}



\begin{quote}
Before the new one: say the Kannada for stubborn, then the Kannada for rotten.
\end{quote}

\subsection*{You'll want to know: \kn{ಮಾಗಿದ}}
\textbf{\kn{ಮಾಗಿದ}} — \emph{māgida} — \textquotedblleft{}ripe\textquotedblright{}.

\begin{grammarlens}[title={What you've built}]
One more word for this chapter.
\end{grammarlens}

\subsection*{Your turn}
\begin{itemize}
\raggedright
  \item \emph{Say it:} \emph{māgida}
  \item \emph{Say it:} \emph{māgida}, once more
  \item \emph{Say it:} say \emph{māgida}
  \item \emph{Recall:} say the Kannada for stubborn, then the Kannada for rotten, then say \emph{māgida} again
\end{itemize}

\begin{tcolorbox}[breakable,colback=teal!4,colframe=teal!35!black,title={Before you move on}]
What does \kn{ಮಾಗಿದ} mean? (\textquotedblleft{}Ripe\textquotedblright{}.) Say it once more.
\end{tcolorbox}
\section[\texorpdfstring{\kn{ಸ್ವಂತ} (svanta)}{ಸ್ವಂತ (svanta)}]{\kn{ಸ್ವಂತ} (svanta) — one's own}

\label{lesson:KA-C182-svanta}



\begin{quote}
Before the new one: say the Kannada for rotten, then the Kannada for ripe.
\end{quote}

\subsection*{You'll want to know: \kn{ಸ್ವಂತ}}
\textbf{\kn{ಸ್ವಂತ}} — \emph{svanta} — \textquotedblleft{}one's own\textquotedblright{}.

\begin{grammarlens}[title={What you've built}]
One more word for this chapter.
\end{grammarlens}

\subsection*{Your turn}
\begin{itemize}
\raggedright
  \item \emph{Say it:} \emph{svanta}
  \item \emph{Say it:} \emph{svanta}, once more
  \item \emph{Say it:} say \emph{svanta}
  \item \emph{Recall:} say the Kannada for rotten, then the Kannada for ripe, then say \emph{svanta} again
\end{itemize}

\begin{tcolorbox}[breakable,colback=teal!4,colframe=teal!35!black,title={Before you move on}]
What does \kn{ಸ್ವಂತ} mean? (\textquotedblleft{}One's own\textquotedblright{}.) Say it once more.
\end{tcolorbox}
